The rapid progress in the field of aerial robotics has enabled their successful utilization in diverse applications including search-and-rescue~\cite{delmerico2019jfr,tranzatto2022cerberus}, inspection~\cite{SIP_AURO_2015,caprari2012highly}, and forest monitoring~\cite{Zhang2016Forest,aucone2023dronedna}. Keys to this success have been techniques to plan paths, often via sampling-based methods and motion primitives~\cite{RRT,karaman2011rrt_star,paranjape2015motion}, within dense $3\textrm{D}$ maps~\cite{voxblox, hornung13auro, museth2013vdb,florence2018nanomap}. Yet, despite the progress, most approaches are limited by a) the need for consistent localization and mapping, b) the latency such systems introduce, as well as c) the maximum resolution of the map used for collision checking and the computational burden of increasing such resolution in order for narrow cross-section (thin) obstacles (e.g., wires, railings, rods, tree branches) to be represented. Recently, we have seen great advances coming from the field of deep learning for collision-free flight tailored to a host of demanding tasks including drone racing~\cite{Foehn2021AlphaDrone}. Yet, most such methods also assume consistent pose estimation~\cite{Loquercio2021Science}, and often rely fully on simulated data for training~\cite{ORACLE_ICRA2022} which, however, may fail to capture significant imperfections of range sensing such as holes in depth maps~\cite{yang2018segstereo,ahn2019analysis}, while there is a lack of focus on methods that explicitly tackle avoiding thin, hard-to-perceive, obstacles.



%%%%%%%%%%%%%%%%%%%%%%%%%%%%%%%%%%%%%%%%%%%%%%%%%%%%%%%%%%%%%%%%%
\begin{figure}[t!]
\includegraphics[width=\columnwidth]{figures/IntroFig_v9_corrected.png}
\vspace{-5ex}
\caption{Instance of an experiment demonstrating collision-free flight in a cluttered environment further involving thin obstacles (e.g., vertical wooden sticks) with the robot only exploiting real-time depth maps and estimates of its linear and angular velocities, as well as attitude.}
\label{fig:intro}
\vspace{-6ex}
\end{figure}

%%%%%%%%%%%%%%%%%%%%%%%%%%%%%%%%%%%%%%%%%%%%%%%%%%%%%%%%%%%%%%%%%

Motivated by the above, this work contributes a modularized learning-based navigation method that a) is capable of navigating cluttered environments involving hard-to-perceive thin obstacles, while b) not requiring access to a (online or offline) map of the environment or the estimation of the full pose of the robot but instead relies solely on real-time depth observations and a partial estimate of the robot's linear/angular velocities and attitude. While the community has considered the explicit detection of certain thin obstacles (e.g., wires are studied in~\cite{Madaan2017WireDU,Kulkarni2020} using RGB images), the challenge of safe navigation in cluttered scenes with complex topologies and morphologies of thin obstacles is not solved. Furthermore, for such structures, acquiring accurate depth maps using real sensors is particularly difficult due to their narrow cross-sections. To overcome this limitation and considering the use of depth cameras, we first create and aggregate both real and simulated datasets of depth frames with pixel-wise labels on thin obstacles. Then, we contribute a semantically-enhanced Variational Autoencoder (seVAE) that a) learns to compress, encode and reconstruct depth maps while preserving thin obstacle features, b) utilizes both real and simulated depth data for training along with pixel-wise labels reflecting thin obstacles (when available), and c) can work directly with raw depth frames while remaining robust to complex cases of sensor noise and errors including holes in the depth maps. Through the encoder of the seVAE, the method arrives --at inference time-- to a low-dimensionality latent space that encodes the information for collision avoidance including against thin obstacles. This latent space is then combined with a Collision Prediction Network (CPN) building upon our previous work called ``ORACLE''~\cite{ORACLE_ICRA2022}. The CPN utilizes the compressed latent space of the encoder combined with the robot's partial state (considering its uncertainty) and provides collision scores for candidate $3\textrm{D}$ action sequences from a motion library thus enabling safe autonomous flight. 


To validate the proposed approach, a set of evaluation studies, both in simulation and experimentally, were conducted. In simulation, the method was tested in environments with increasing density of obstacles and its results are compared against an end-to-end trained $3\textrm{D}$ extension of ORACLE (using simulated depth images only). Two real-world experiments were also performed to test the new method in densely cluttered environments with various sizes and types of obstacles. Similarly, the performance of the new method in these experiments is compared against that of the end-to-end trained ORACLE demonstrating the value added by the ability to incorporate real data in the training phase to enhance robustness against sensor imperfections.



In the remainder of this paper, Section~\ref{sec:related} presents related work, followed by the proposed modularized learning-based navigation solution in Section~\ref{sec:approach}. Evaluation studies are detailed in Section~\ref{sec:evaluation}, and conclusions in Section~\ref{sec:conclusions}.

